\section{读懂配合物}
\learning{能从化学式辨认内外界，分别计算氧化态、配位数与 $d$ 电子数；能把配体的连接方式读出来，并完成简单配合物的双向命名。}
\sources{A，第12--19、26--32、37--46、58--78页；B，第1--3页。}

\subsection{先分清：一个配合物里有哪些结构单元}
配合物由中心原子与配体通过配位作用结合而成。中心通常是金属原子或离子；配体可以是中性分子，也可以是离子。方括号内的\textbf{配位个体}通常称为内界；为平衡电荷而存在、没有直接与该中心配位的离子称为外界。例如，\ce{[Co(NH3)5Cl]Cl2} 的内界是 \ce{[Co(NH3)5Cl]^2+}，外界有两个 \ce{Cl-}。内界本身也可以中性，如 \ce{[Ni(CO)4]}，此时不需要外界反离子。

三个容易混用的计数要分开。\textbf{配位原子}是配体上直接与金属结合的原子；\textbf{齿数}是一个配体实际向同一个金属提供的配位原子数；\textbf{配位数}则是该金属周围直接相连的配位原子总数。乙二胺 en 以两个 N 配位，是双齿配体；\ce{[Co(en)3]^3+} 含三个配体，却是六配位。桥联配体连接多个金属时，要对每个金属分别计数，不能把整个团簇共用的原子只算一次。

螯合配体用两个或更多配位原子结合\emph{同一个}金属，并形成包含金属的环。两可配体则可以换用不同原子配位，如 \ce{NO2-} 的 N 配位与 O 配位；“可以换一个原子”不等于“同时以两个原子螯合”。en 常形成五元螯环；五元、六元环通常较容易适应配位几何，但环大小不是稳定性的唯一判据，还要考虑键角、柔性、金属大小及配体的预组织。

\subsection{氧化态与电子数：先算电荷，再谈成键}
给定配位个体的总电荷 $q$，按约定给配体分配形式电荷 $q_i$，中心金属的氧化态 $x$ 满足
\[
 x+\sum_i q_i=q.
\]
对常见过渡金属，形式 $d$ 电子数可由“族序数减去氧化态”得到。例如，Co 属第9族，Co(III) 为 $d^6$；Ni 属第10族，Ni(II) 为 $d^8$，Ni(0) 为 $d^{10}$。\ce{NH3}、水、en、CO 的形式电荷为0，卤离子为 $-1$，草酸根 ox 为 $-2$。不能因为 CO 能接受反馈电子，就把它按带负电的配体来配平。

氧化态是电子归属的\textbf{记账规则}，不是金属真实电荷的测量值。含配体自由基等非无辜配体的体系，可能存在“金属被氧化”与“配体被氧化”的不同描述；此时需要结合光谱和电子结构证据，不能只靠总电荷认定电子实际在哪里。

\begin{worked}[同一式子，回答四种问题]
求 \ce{[CoCl2(en)2]Cl} 的内外界、Co 氧化态、配位数和 $d$ 电子数。

外界为一个 \ce{Cl-}，所以内界 \ce{[CoCl2(en)2]+} 带 $+1$ 电荷。由 $x+2(-1)+2(0)=+1$ 得 Co 为 $+3$ 价。两个 en 各提供两个 N，两个氯各提供一个 Cl，因此配位数是 $2\times2+2=6$，而配体数是4。Co(III) 的形式电子组态为 $d^6$。式子尚未写出 cis 或 trans，因而\emph{没有唯一指定几何异构体}。
\end{worked}

\subsection{常见配体：记住电荷、给体原子与连接限制}
下表将常见的形式电荷、给体原子和连接约束放在一起；实际齿数仍由具体结构确定。

\begin{center}
\small
\begin{tabularx}{\linewidth}{@{}l c l X@{}}
\toprule
配体或缩写 & 电荷 & 常见给体／齿数 & 做题时要留意的事 \\
\midrule
\ce{NH3}、\ce{H2O} & 0 & N、O；单齿 & 氨与水均不带电 \\
\ce{Cl-}、\ce{Br-} & $-1$ & 卤素；单齿 & 可端基配位，也可桥联 \\
\ce{CN-} & $-1$ & 常为 C；单齿 & 也可经 N 配位或桥联 \\
\ce{NO2-} & $-1$ & N 或 O；常为单齿 & 需标明是 N 配位还是 O 配位 \\
\ce{SCN-} & $-1$ & S 或 N；常为单齿 & SCN 与 NCS 的写法提示连接端 \\
CO & 0 & C；单齿 & $\sigma$ 给电子、$\pi$ 接受电子 \\
en & 0 & N,N；双齿 & 乙二胺，柔性五元螯环 \\
bpy、phen & 0 & N,N；双齿 & 2,2$'$-联吡啶、1,10-菲咯啉，骨架较刚性 \\
ox、acac & $-2$、$-1$ & O,O；双齿 & 分别为草酸根、乙酰丙酮根 \\
terpy & 0 & N,N,N；三齿 & 三联吡啶，通常要求近似 mer 的排列 \\
trien、tren & 0 & 四个 N；四齿 & 线性链与支化骨架，几何约束不同 \\
edta & $-4$ & 常为 N$_2$O$_4$；六齿 & 必须区分 edta$^{4-}$ 与其质子化形式 \\
\bottomrule
\end{tabularx}
\end{center}

\subsection{配位键不只是“画一个箭头”}
用 $L\to M$ 表示配体向金属提供电子对，强调的是成键轨道的来源。成键以后，电子处在整个配合物的分子轨道中；配位键并不是一种永远只能由一方“拥有电子”的特殊共价键。实际金属--配体键通常同时具有共价性和静电作用，也不能仅凭某个固定的电负性差阈值把所有键二分。

氨可用 N 上的孤对电子向金属提供 $\sigma$ 相互作用。CO 则既以 C 端孤对电子给电子，又用空的 $\pi^*$ 轨道接受金属 $d_\pi$ 电子，形成\textbf{$\pi$ 反馈}。两种作用常相互协同；反馈会占据配体的反键轨道，使 C--O 键趋于减弱。烯烃配位也是类似的授受组合：烯烃的已占据 $\pi$ 轨道向金属给电子，金属向烯烃空的 $\pi^*$ 轨道反馈。蔡司盐中的乙烯正是这一模型的代表。这里“$\sigma$ 给电子”描述相对于金属--烯烃连接方向的对称性，并不意味着电子必须来自烯烃的 $\sigma$ 键。

卤离子等配体还可通过已占据的 $\pi$ 轨道向金属给电子，属于 $\pi$ 给体；CO 等接受金属反馈电子的配体则属于 $\pi$ 受体。比较配体时，应分别说明 $\sigma$ 与 $\pi$ 相互作用中的电子流向。

\textbf{配位层}帮助我们从局部结构逐层看向整个配体。第一层关注金属与直接配位原子；第二层关注把这些原子连接起来的螯合骨架；更外层包括取代基等部分，它们可以改变位阻、诱导效应及分子间作用。\ce{[Co(NH3)6]^3+} 与 \ce{[Co(en)3]^3+} 都有 CoN$_6$ 第一层，但后者的 N 被成对连起来，能形成有手性的螺旋排列。这种结构差异要到第二层才看得出来。

\subsection{Werner 的证据链：从组成走向空间结构}
Werner 在1893年提出配位理论，把当时的主价、副价分别与今日的氧化态、配位数相联系，并强调配位有确定的空间方向。其说服力不只来自某一种颜色或某一次分析，而来自组成、反应、溶液性质与异构体数目之间的相互制约。

\begin{center}
\small
\begin{tabularx}{\linewidth}{@{}l c c X@{}}
\toprule
化合物 & 迅速沉淀的 Cl$^-$ & 理想离子数 & 结构解释 \\
\midrule
\ce{[Co(NH3)6]Cl3} & 3 & 4 & 三个氯均在外界 \\
\ce{[CoCl(NH3)5]Cl2} & 2 & 3 & 一个氯进入内界，仍为六配位 \\
\ce{[CoCl2(NH3)4]Cl} & 1 & 2 & 两个氯在内界，可出现 cis/trans \\
\ce{[CoCl3(NH3)3]} & 0 & 0 & 中性内界，可出现 fac/mer \\
\bottomrule
\end{tabularx}
\end{center}

表中的沉淀量指每摩尔化合物在适当条件下与 \ce{Ag+} 快速反应所得 \ce{AgCl} 的摩尔数；离子数按内界保持完整、理想电离计算。电导率与依数性支持不同的外界离子数，但电导率并不与离子数严格成比例。六配位系列的几何异构体数进一步支持八面体，而不是只凭“六配位”就宣布必为八面体。1911年对手性钴配合物的成功拆分，又证明配体确实具有可分辨的三维排列；后来不含碳的配合物也被拆分，说明手性不以不对称碳为必要条件。

\subsection{从式到名，再从名还原化学式}
命名可分成两层：先说清\emph{一个配位个体}中的配体、数目、连接方式和金属氧化态，再按盐的组成连接内外界名称。本笔记沿用讲义的中文配体排列习惯：无机配体在有机配体之前，同类中通常先阴离子后中性配体，再依据给体原子等规则排列。不要把这套中文顺序与英文系统命名的字母排序拼接使用。复杂配体的数目用括号隔开，例如“三(乙二胺)”，避免把配体名称中的数字误读为个数。

简单外界阴离子常用“氯化……”这样的盐名，含氧酸根用“硫酸……”“硝酸……”；配阴离子的盐通常写作“……合某金属酸某”。例如 \ce{K2[PtCl6]} 为六氯合铂(IV)酸钾，\ce{[Co(NH3)6]Cl3} 为氯化六氨合钴(III)。几何信息用 cis、trans、fac、mer 放在名称或配位个体之前，不能指望只写组成就能还原空间排列。

连接标记中，\textbf{$\kappa$ 指出具体的配位原子}，如亚硝酸根-$\kappa N$ 与亚硝酸根-$\kappa O$；en 的 $\kappa^2N,N'$ 表示两个 N 都与同一中心结合。\textbf{$\eta$ 表示同一连续片段参与配位的原子数}，如二茂铁 \ce{[Fe(C5H5)2]} 中每个茂环为 $\eta^5$，乙烯常为 $\eta^2$。它不是两个孤立给体原子的齿数记号。讨论 $\pi$ 配体的“配位数”时应先声明按原子还是按配位位点计数；不能把传统的空间位点数与逐原子计数混加。

\textbf{$\mu$ 表示桥联}，$\mu_2$ 表示连接两个中心，通常可简写为 $\mu$。例如 \ce{[(NH3)5Cr(OH)Cr(NH3)5]Cl5} 可称五氯化($\mu$-羟)·二[五氨合铬(III)]；一个桥联 \ce{OH-} 同时为两个 Cr 各贡献一个 O 配位原子。名称中的“羟”不应被误读为中性水。

\begin{worked}[把名字中的信息逐项还原]
“氯化氯·五氨合钴(III)”对应什么式子？“反-四(异硫氰酸根)·二氨合铬(III)酸铵”又怎样还原？

第一例的内界含 Co(III)、一个 \ce{Cl-} 和五个中性氨，净电荷为 $+2$，需两个外界氯离子，故为 \ce{[CoCl(NH3)5]Cl2}。开头的“氯化”与内界名称中的“氯”不是重复计数。

第二例内界电荷为 $+3-4=-1$，故为 trans-\ce{NH4[Cr(NCS)4(NH3)2]}。“异硫氰酸根”在此表示经 N 配位，也可用 $\kappa N$ 明示；trans 指两个氨互处反位。检验时同时检查总电荷为零、配位数为6和 N 配位方式，三项都满足才算还原完整。
\end{worked}

\section{从配位数到空间构型}
\learning{知道常见配位数对应哪些候选构型，能解释四配位和六配位的主要选择，并把金属电子结构、配体约束与结构互变联系起来。}
\sources{B，第1--9、27--32页；A，第73--78页。}

\subsection{配位数给候选，不给唯一答案}
确定空间构型时，先数金属周围的实际配位原子，再比较候选多面体。下面的名称说的是\emph{配位原子的位置}，不是整个配合物的外形；实际结构允许有键长、键角和扭转的偏离。

\begin{center}
\small
\begin{tabularx}{\linewidth}{@{}c l X@{}}
\toprule
配位数 & 主要候选构型 & 如何识别或理解 \\
\midrule
2 & 直线形 & 两个方向近似相反，Au(I)、Ag(I) 常见 \\
3 & 平面三角形；三角锥 & 前者金属与三个给体近共面，不是四面体少画一键 \\
4 & 四面体；平面四方形 & 是否共面、是否有 $180^\circ$ 反位关系是关键 \\
5 & 三角双锥；四方锥 & 分别为三赤道加两轴向、四基面加一顶点 \\
6 & 八面体；三棱柱 & 前者有三对反位，后者有两个近重叠的平行三角面 \\
7 & 五角双锥；单帽八面体；单帽三棱柱 & “加帽”是在原多面体某一面的外侧增加顶点 \\
8 & 四方反棱柱；三角十二面体等 & 也可见立方体、双帽三棱柱、六角双锥 \\
\bottomrule
\end{tabularx}
\end{center}

\subsection{四配位：为什么 $d^8$ 不是充分条件}
四面体的四个配位方向彼此等价，没有一对配体严格处于反位；平面四方形则有两对相对位置，因而能支持常见的顺反异构。几何选择要比较电子能量与配体间排斥，不能仅凭配位数区分，也不能把杂化轨道标签当成结构选择的独立原因。

对 $d^8$ 金属，平面四方形分裂常把一个强反键的 $d$ 轨道抬得很高，使八个电子占据其余轨道而获得有利排布。Pd(II)、Pt(II) 的四配位配合物通常倾向平面四方形；Ni(II) 则对配体及环境更敏感。典型地，\ce{[Ni(CN)4]^2-} 为平面四方形、抗磁，\ce{[NiCl4]^2-} 为四面体、顺磁：二者都是 Ni(II)、$d^8$，电子数并没有改变。

在一些 Ni(II) 配合物中，平面形与四面体形的能量差不大，改变配体、溶剂或温度便可能改变平衡。讲义中的膦配合物体现了这种构型与磁性的联系。磁性可以作为构型线索，再与组成、光谱及衍射证据相互印证。

\subsection{五配位互变、六配位的两个三角面}
三角双锥中，轴向配体与三个赤道配体分别约成 $90^\circ$，赤道配体之间约成 $120^\circ$；四方锥则有四个近共面的基面位置与一个顶端位置。真实五配位结构常位于两者之间。\textbf{Berry 假旋转}使两个轴向位置与两个赤道位置交换，途经近似四方锥的几何，另一个赤道位置在这次过程中不参与交换。它不需要先把配体完全解离；若互变比谱学观测快，原本不等价的位置可能表现为时间平均等价。

六配位以八面体最常见，但不是只有八面体。可把八面体看作两个相对的三角面，沿共同轴投影时呈交错排列；三棱柱的两个三角面近似重叠。两者之间的扭转提供了理解构型变形的直观路径。讲义中的三(二硫烯)铼配合物 \ce{[Re(S2C2(CF3)2)3]} 是三棱柱结构的例子。某些金属与受限的螯合骨架会稳定这种排列，所以看到六个配位原子时，仍要检验反位关系和三角面的相对取向。

\subsection{七、八配位：用“底面、扭转、加帽”读图}
七配位可以从一个五边形加上下两个顶点得到五角双锥；也可以从六配位多面体出发，在一个面外加帽。同样是七配位，晶格与反离子环境也可能改变哪一种结构更稳定。配位数相同，不保证两个晶体拥有同一多面体。

八配位的\textbf{四方反棱柱}可由立方体帮助想象：保留上下两个正方形，把其中一个绕中心轴转动约 $45^\circ$，上下顶点便交错排列；连接相邻顶点后，侧面由四边形变成三角形。\textbf{三角十二面体}有八个顶点、十二个三角面，名称中的“十二”指面数，不是配位数。可把立方体的顶点交替分成两组四面体，再对两组作不同的拉伸、压缩来理解其变形。它也不是有二十个顶点的正十二面体。双帽三棱柱是在三棱柱基础上再加两个帽；六角双锥则是六个赤道点加两个轴向点。

\subsection{结构选择时按什么顺序分析}
先看\textbf{金属容纳空间}：较大的金属与较小的给体往往较容易形成高配位数；氧化态升高可加强吸引，却也会改变金属半径，不能只记“电荷越高，配位数越大”。稀土和较大的早过渡金属常见高配位数。$d^0$ 体系没有传统的 $d$ 电子配体场稳定化能，不能把它们的高配位直接归因于“得到更多配体场稳定化能”。

再看\textbf{配体是否装得下、够不够得着}：体积大的配体会增加排斥，多齿配体的齿间距离和咬合角则限制可占据的位置。最后比较\textbf{电子结构与环境}：$d$ 电子数、$\sigma/\pi$ 成键、溶剂配位、反离子及晶体堆积都会参与选择。

\section{枚举异构体与判断手性}
\learning{先按连接关系分类，再按几何关系与镜像关系计数；会用明确坐标枚举八面体异构体，并把单分子立体异构体与样品混合物分开。}
\sources{B，第10--25、29--48、61--65页；A，第37--44、67--78页。}

\subsection{第一道分岔：连接关系变了没有}
同分异构体具有相同的总组成，却不是同一个结构。先问金属与哪些原子相连、哪些物种在内界；若这些连接不同，属于结构异构。若连接相同而空间排列不同，才进入立体异构的讨论。

\begin{center}
\small
\begin{tabularx}{\linewidth}{@{}l X X@{}}
\toprule
类型 & 一组典型例子 & 实际改变了什么 \\
\midrule
配位异构 & \ce{[Co(NH3)6][Cr(ox)3]} 与 \ce{[Cr(NH3)6][Co(ox)3]} & 配阳、配阴离子之间交换了配体的归属 \\
键合异构 & \ce{[Co(NO2)(NH3)5]^2+} 与 \ce{[Co(ONO)(NH3)5]^2+} & 同一个两可配体换用 N 或 O 配位 \\
电离异构 & \ce{[CoBr(NH3)5]SO4} 与 \ce{[Co(SO4)(NH3)5]Br} & 内外界交换，直接释放的外界离子不同 \\
水合异构 & \ce{[Cr(H2O)6]Cl3} 与 \ce{[CrCl(H2O)5]Cl2.H2O} & 水在配位水与结晶水之间改变归属 \\
\bottomrule
\end{tabularx}
\end{center}

水合异构属于溶剂合异构的特例。比较时应把结晶水计入总组成，不能因为它在方括号外就忽略。配体自身若换成另一种连接异构体，例如1,2-丙二胺换成1,3-丙二胺，相应配合物也改变了结构；这与同一配体换配位原子的键合异构不是同一层次。

\subsection{把 cis/trans 和 fac/mer 放在明确的位置上}
平面四方形中，cis 表示指定的两个配体相邻，夹角约 $90^\circ$；trans 表示相对，约 $180^\circ$。四面体中任意两个顶点的关系相同，不能照搬平面形的 cis/trans。下面两图只画金属与配位原子的连接，不表示键长比例。

\begin{center}
\begin{tikzpicture}[every node/.style={font=\small},line width=.6pt]
\begin{scope}[xshift=-2.6cm]
\node (m) at (0,0) {Pt};
\node (t) at (0,1) {Cl}; \node (r) at (1.3,0) {Cl};
\node (b) at (0,-1) {$\mathrm{NH_3}$}; \node (l) at (-1.3,0) {$\mathrm{H_3N}$};
\draw (m)--(t) (m)--(r) (m)--(b) (m)--(l);
\node at (0,-1.65) {cis：两个 Cl 相邻};
\end{scope}
\begin{scope}[xshift=2.6cm]
\node (m) at (0,0) {Pt};
\node (t) at (0,1) {Cl}; \node (r) at (1.3,0) {$\mathrm{NH_3}$};
\node (b) at (0,-1) {Cl}; \node (l) at (-1.3,0) {$\mathrm{H_3N}$};
\draw (m)--(t) (m)--(r) (m)--(b) (m)--(l);
\node at (0,-1.65) {trans：两个 Cl 相对};
\end{scope}
\end{tikzpicture}
\end{center}

八面体可统一用六个位置 $+x,-x,+y,-y,+z,-z$ 表示：同一轴的正负端互为 trans，不同轴的任意两端互为 cis。对 $MA_3B_3$，三个 A 占据一个三角面时为 fac；三个 A 中有一对反位、三者与金属共处一个经面时为 mer。下面的占位表可直接用来画模型，未列出的配体填满其余位置。

\begin{center}
\small
\begin{tabular}{@{}lll@{}}
\toprule
组成与形式 & 指定配体的占位 & 一眼检查的条件 \\
\midrule
$MA_4B_2$，trans & B 在 $+z,-z$ & 两个 B 相对 \\
$MA_4B_2$，cis & B 在 $+z,+x$ & 两个 B 相邻 \\
$MA_3B_3$，fac & A 在 $+x,+y,+z$ & A 之间三组关系全为 cis \\
$MA_3B_3$，mer & A 在 $+x,-x,+z$ & A 中恰有一组 trans \\
\bottomrule
\end{tabular}
\end{center}

\subsection{计数先固定模型，再消去重复}
枚举时先说明配体是否单齿、双齿是否对称，是否忽略螯环构象与配位 N 的手性。随后固定一个配体位置，依据剩余位置的相邻、相对关系分类，并用\emph{整体旋转}消去重复。镜面反射不能用于把两个对映体合并；只有在专门统计“几何类型”时，才把一对对映体归入同一类。

下表把几何类型与分子立体异构体数并列。A、B 等为不同的无手性单齿配体，AA 为对称双齿配体；忽略环构象和配体自身手性。“几何类型”一栏按不分镜像的一对计一个类型。

\begin{center}
\small
\begin{tabularx}{\linewidth}{@{}l c c X@{}}
\toprule
模型 & 几何类型 & 立体异构体 & 计数原因 \\
\midrule
平面 $MA_2B_2$ & 2 & 2 & cis、trans，均无手性 \\
平面 $MABCD$ & 3 & 3 & 固定 A，再选其反位配体，有三种选择 \\
四面体 $MABCD$ & 1 & 2 & 一对不能重合的镜像 \\
八面体 $MA_4B_2$ & 2 & 2 & cis、trans，均无手性 \\
八面体 $MA_3B_3$ & 2 & 2 & fac、mer，均无手性 \\
八面体 $MA_2B_2C_2$ & 5 & 6 & 全反1种、恰一对反位3种、全顺1对对映体 \\
八面体 $M(AA)_3$ & 1 & 2 & $\Delta$、$\Lambda$ 一对 \\
八面体 $M(AA)_2B_2$ & 2 & 3 & trans 1个；cis 为一对对映体 \\
八面体 $MABCDEF$ & 15 & 30 & 六种配体全不同，15对对映体 \\
\bottomrule
\end{tabularx}
\end{center}

对最后一行，可以先把六个配体分为三组反位配对，共有 $5\times3\times1=15$ 种。每个配对方案仍有两个相反的空间手性；不同配体全部标记后，不存在能把两者重合的镜面或反心。因此分子总数是30，不是15。这个思路比画30个没有位置标记的八面体更容易查错。

\subsection{镜像、手性与构象：三个问题不要互相代替}
\textbf{对映体}互为镜像而不能经整体旋转重合；\textbf{非对映异构体}是彼此不呈镜像关系的立体异构体，例如 cis 与 trans。非对映异构体不要求双方都有手性。判断一个指定结构是否有手性，应检查它能否与镜像重合；等价的对称性判据是不存在任何反轴 $S_n$，其中也包括镜面 $S_1$ 和反心 $S_2$。只有“没有反心”远远不够，含镜面的结构仍然无手性。

无手性配体也能排出有手性的整体。\ce{[Co(en)3]^3+} 的三个螯环像螺旋桨，形成右手 $\Delta$ 与左手 $\Lambda$ 两种构型。判读时沿选定的三重轴看螯合带向远处延伸的旋向，并与右手或左手螺旋比较；先标明观察方向，不能只凭纸面上线条的斜率猜。大写 $\Delta/\Lambda$ 描述整体螯合排列，小写 $\delta/\lambda$ 可描述五元螯环的局部构象，二者不是同一个标签。对于 edta 等多齿配体，还需要按明确的螯环组合规则指定，不能随意套用三条等价螯合带的图像。

\textbf{构象异构}来自单键转动或环翻转等变化，通常不改变原子连接，也不一定改变金属中心构型。例如 en 环的翻转可以改变环褶皱而不改变整个三螯合体的 $\Delta/\Lambda$。若转动能垒很高，阻转异构体可在实验时间尺度上分离；若互变很快，谱学看到的可能只是平均结构。是否能分离，必须连同温度与观测时间说明。

\begin{keyidea}[构型符号不是旋光方向]
$\Delta/\Lambda$、$R/S$ 是空间构型标签；$(+)/(−)$ 是在指定波长、溶剂、温度等条件下测得的旋光符号，二者没有普遍的一一对应关系。同一构型在不同波长下甚至可能改变旋光符号。外消旋样品含等量对映体，净旋光可相消，但其中的每一个分子仍可有手性。$rac$ 是样品组成说明，不是新增的分子异构体。
\end{keyidea}

\subsection{trien 与 tren：同为四齿，限制却不同}
trien 是线性的三乙四胺，四个 N 按 $N_1$--$N_2$--$N_3$--$N_4$ 连成一条链；tren 是三(2-氨乙基)胺，一个中心 N 连着三条含端基 N 的支臂。二者都能提供四个 N，但不能任意交换它们在八面体上的位置。下面采用讲义的常见五元螯环模型，不计配位 N 的独立手性，也不把不同环褶皱拆成新异构体。

\begin{worked}[{用坐标完整枚举 $[MX_2(\mathrm{trien})]$ 与 $[MX_2(\mathrm{tren})]$}]
以 $\pm x,\pm y,\pm z$ 为八面体六个顶点，X 是两个相同的单齿配体。trien 的相邻 N 必须占据相邻顶点，以便形成三个常见五元螯环。以下每行给出一份可重建的代表结构；连线顺序严格是 $N_1\to N_2\to N_3\to N_4$，而不是任意把四个 N 两两相连。

\begin{center}
\footnotesize
\begin{tabular}{@{}llll@{}}
\toprule
形式 & 两个 X 的位置 & trien 四个 N 的链路径 & 理想点群 \\
\midrule
trans & $+z,-z$ & $+x\to+y\to-x\to-y$ & $C_{2v}$ \\
cis-a & $+x,+z$ & $+y\to-x\to-z\to-y$ & $C_2$ \\
cis-b & $+x,+z$ & $+y\to-x\to-y\to-z$ & $C_1$ \\
\bottomrule
\end{tabular}
\end{center}

cis-a 的两个末端 N 互处反位，cis-b 的两个末端 N 互为顺位，因此两者不能靠整体旋转变成同一种连接排列。trans 含镜面，无手性；两种 cis 均有手性。将坐标作镜像变换 $(x,y,z)\mapsto(x,-y,z)$，即可分别得到 cis-a 与 cis-b 的另一对映体。故 trien 有 \textbf{3种几何类型，5个分子立体异构体}：$1+2+2$。

对 tren，中心 N 必须与三个支臂末端 N 分别形成螯环。可取中心 N 在 $+z$，三个末端 N 在 $+x,-x,+y$，剩余两个 X 在 $-y,-z$，互为 cis。三个相同支臂的置换不会产生新的几何类型；在此模型下理想点群为 $C_s$，有镜面，故只有\textbf{1种几何类型、1个无手性分子立体异构体}。若进一步计入柔性环的冻结构象，可能出现额外的构象手性，但那是另一层计数。

\end{worked}

\subsection{如何用实验证据区分异构体}
衍射给出连接与空间排列；溶液 NMR 比较等价环境及其相互关系，快速交换时需采用时间平均结构解释。偶极矩利用键偶极的矢量相消区分某些顺反异构体；IR 则比较目标振动的对称性与谱带数，具体计算见第9节。电子光谱提供能级与跃迁强度的信息，可与上述证据配合使用。成功拆分能证明存在可分离的对映体，拆分失败则不足以判断有无手性。

\section{手性表征与不对称催化}
\learning{分清分子手性、晶体手性与样品对映组成，知道绝对构型的证据来源；会计算 ee 与 dr，并能解释普通动力学拆分的转化率限制。}
\sources{B，第33--39、49--60、66--70页；A，第38--45、52页。}

\subsection{外消旋样品结晶后，有三种不同的固态组织}
设一对对映体为 R 与 S。总量各半只规定总体组成，不规定它们怎样进入晶体。

\begin{center}
\small
\begin{tabularx}{\linewidth}{@{}l X X@{}}
\toprule
固态形式 & 单颗晶体内部 & 对分离的含义 \\
\midrule
外消旋混合物（聚集体） & 理想情况下，R 晶体与 S 晶体各自形成 & 可有自发拆分；整个晶体批次仍可为1:1 \\
外消旋化合物 & R、S 在同一晶格中按确定比例有序排列 & 不能靠挑出两类对映纯晶体直接拆分 \\
外消旋固体溶液（假外消旋体） & R、S 可统计性占据晶格位置，组成可变 & 不能仅凭单颗晶体就假定其对映纯 \\
\bottomrule
\end{tabularx}
\end{center}

自发拆分说的是对映体在结晶时分别成晶，并不自动等于整批产物获得非零 ee。若两类晶体恰好等量，混合溶解后仍不旋光。机械挑晶、优先结晶等方法的适用性依赖固相性质；另一种常见拆分思路是加入对映纯拆分剂，把一对对映体转成性质不同的非对映盐或配合物，分离后再除去拆分剂。Werner 钴配合物的经典拆分就利用了这种差别。

还要区别\textbf{分子手性与晶体手性}。无手性的分子或离子也可能按手性方式堆积；讲义中的氯酸钠晶体就是代表。晶体溶解后若手性仅来自堆积，堆积被破坏，溶液不必保留旋光。反之，对映纯且构型稳定的手性配合物溶解后仍可保持分子手性。只含平移、旋转、螺旋操作的 Sohncke 空间群允许手性晶体，但空间群名称本身不能证明分子对映纯，更不能替代绝对构型测定。

\subsection{绝对构型的三类证据，以及 HPLC 能回答什么}
\textbf{反常 X 射线衍射}利用反常散射造成的衍射强度差别区分相反的绝对结构，是直接测定的重要方法。可靠判定需要足够的反常信号、合适的数据与结构模型，并通过绝对结构参数等指标评估；元素、波长、晶体质量和孪晶情况都会影响结果。

\textbf{电子圆二色 ECD／旋光色散 ORD}可通过已知参照物的 Cotton 效应或理论计算间接关联构型。参照物应有相近的空间结构、电子结构与跃迁归属。\textbf{振动圆二色 VCD}比较左、右圆偏振红外吸收差，常结合构象搜索和计算谱，判定哪一绝对构型更符合实验。这两类方法都需要把样品的主要构象与溶剂条件纳入比较。

\textbf{手性 HPLC}用于分离对映体并测定比例；绝对构型则需要通过已知标准品、可靠关联或联用证据指认，不能只根据先后出峰命名。测定比例时应保证分离充分，并确认检测响应。

\subsection{ee 与 dr：先说明分母里有哪些物种}
一对对映体的对映体过量定义为
\[
 ee=\frac{|n_R-n_S|}{n_R+n_S}\times100\%.
\]
报告时应同时指出哪一个占优；也可用带符号形式表达方向，但要说明约定。若 R、S 是唯一两种组分，则主对映体与次对映体的摩尔分数分别为 $(1+ee)/2$ 和 $(1-ee)/2$，此处 ee 代入小数形式。ee 不是主对映体含量，也不是产率。

非对映比 dr 比较指定非对映异构体的量，例如 syn:anti 或某两个相对构型家族。一个分子不一定要有两个不对称碳才有非对映异构体；顺反关系、金属构型与配体手性的组合都可能产生非对映关系。若每一非对映家族内部又有一对对映体，dr 与 ee 必须分别报告，不能相互替代。

\begin{worked}[同一组数据，同时求 dr 与各自的 ee]
样品含两组非对映异构体：A 家族的两个对映体为72和8份，B 家族为12和8份；各数字是已经校正响应的相对摩尔量，总量为100份。
\[
 \mathrm{dr}(A:B)=(72+8):(12+8)=80:20=4:1.
\]
A 家族的 $ee_A=(72-8)/80=80\%$，B 家族的 $ee_B=(12-8)/20=20\%$，各以前一个对映体占优。不能把72除以100叫作72\% ee，也不能把四个组分硬合并成一个没有定义的“总 ee”。若只有一对对映体且 $ee=90\%$，其组成才是95:5，而不是90:10。
\end{worked}

\subsection{手性催化怎样让两条反应路径不再等价}
前手性底物在无手性环境中从两个对映面反应，常形成一对对映产物；引入手性金属--配体环境后，两条路径的过渡态成为非对映关系，能垒便可以不同。底物先配位，配体的空间限制、电子作用和定向相互作用共同影响其取向及成键路径。因此，不对称催化并不只是“用大基团挡住一面”，更准确的任务是比较通向两个产物的完整过渡态。

若两条不可逆平行路径分别给出 R、S，且产物生成后不再消旋，速率差可表现为产物比。高 ee 表示选择性高，却不保证转化率、产率或催化剂稳定性高。讲义中的手性铑膦催化不对称氢化说明，金属提供反应活性，手性配体使反应路径不等价；后续结晶还可能进一步富集优势对映体。这些步骤的贡献应分开描述。

\subsection{普通动力学拆分为什么有50\%的单一对映体上限}
不对称催化通常把前手性底物转成手性产物；\textbf{普通动力学拆分}则从外消旋底物出发，使一个对映体反应得快、另一个慢，从而同时富集产物与未反应底物。它并不是“只有一个对映体能反应”的同义词；那只是选择性无限大时的理想极限。动力学拆分也可以使用手性催化剂，所以二者的差别在底物与选择方式，不在于有没有催化剂。

\begin{worked}[转化率超过一半后，纯度为何下降]
从100 mmol 外消旋底物出发，R、S 各50 mmol。设 R 先完全反应，反应保持立体信息且产物不消旋。在50\%转化率时，可得50 mmol 对映纯的 R 来源产物，并留下50 mmol 对映纯 S 底物。

若继续反应至60\%转化率，就必须再消耗10 mmol S。产物现含50 mmol R 来源产物和10 mmol S 来源对映产物，因此
\[
 ee_P=\frac{50-10}{60}\times100\%=66.7\%.
\]
一般地，在这一理想模型的 $c>0.5$ 区间，$ee_P=(1-c)/c$；在 $c<0.5$ 区间，产物仍可对映纯，剩余底物的 $ee_S=c/(1-c)$。两式中的 ee 均为小数。实际选择性有限时，在相同转化率下往往更难兼顾产物与回收底物的对映纯度。
\end{worked}

这里的50\%上限指\emph{不发生对映体互变的普通拆分中，从外消旋起始物取得某一对映纯组分的理论量}，不是所有手性催化反应的产率上限。若反应过程中配合底物消旋，使慢反应对映体持续转为快反应者，则属于动态动力学拆分，物料限制已改变；不能拿这种结果与普通拆分的50\%界限直接比较。
